电流是大量带电粒子整体定向运动的宏观表现。现代生活中的照明、通信、计算和动力设备都依赖电流，但从
物理上理解电流不能只停留在“电荷在导线中流动”这一句话。不同物质中的载流子并不相同：金属主要依靠
电子，电解质由正、负离子共同导电，半导体中电子和空穴都可以形成电流，气体被电离后也会出现电子与离子
电流。载流子的浓度、电荷、质量、碰撞机制以及温度、场强、光照等外部条件都会影响材料的导电性质。

真实导电机制属于量子物理的重要内容。经典模型不能精确解释金属电子的能量分布、晶格散射和能带结构，
但它能够把宏观 Ohm 定律与电子的漂移、碰撞和惯性联系起来，也能清楚说明经典图像在什么条件下失效。
因此本讲同时采用两种层次：用电流、电压和伏安特性研究电路的外部行为，再用简化微观模型解释这些规律。

\subsection*{电流强度与电流密度}

取导电区域中的一个有向截面 $S$。若在时间 $\Delta t$ 内穿过截面的净电荷为 $\Delta Q$，平均电流为
$\Delta Q/\Delta t$；令时间间隔趋于零，得到瞬时电流
\begin{equation}
  \boxed{I(t)=\frac{\dd Q}{\dd t}}.
\end{equation}
电流方向约定为正电荷运动方向。金属中电子带负电，所以电子平均漂移方向与传统电流方向相反。若已知
$I(t)$，从 $t_1$ 到 $t_2$ 穿过截面的净电荷为
\begin{equation}
  \boxed{Q(t_1\to t_2)=\int_{t_1}^{t_2}I(t)\dd t}.
\end{equation}
电流不随时间改变时称为稳恒电流，此时在时间 $t$ 内通过的电量为 $Q=It$。

先假设截面各处载流情况相同，只有一种载流子，数密度为 $n$、单个电荷为 $q$、平均漂移速度为
$\vect v_d$。在时间 $\Delta t$ 内，能穿过面积 $A$ 的载流子来自体积 $Av_d\Delta t$，因此
\begin{equation}
  \boxed{I=nqAv_d}.
\end{equation}
式中的正负号由 $q$ 和速度方向共同决定。它把可测的宏观电流与载流子浓度、单粒子电荷和平均漂移速度联系
起来。

漂移速度通常非常小。以截面积 $A=\SI{1}{mm^2}$ 的银导线为例，若每个银原子近似提供一个导电电子，
电子数密度约为 $n\sim6\times10^{28}\,\si{m^{-3}}$。通入 $I=\SI{1}{A}$ 时
\begin{equation}
  v_d=\frac{I}{neA}\sim10^{-4}\,\si{m.s^{-1}}.
\end{equation}
电子的无规则运动速度却可达 $10^6\,\si{m.s^{-1}}$ 的数量级，远大于漂移速度。因此“有电流”不是说每个
电子都沿导线缓慢而整齐地前进，而是高速无规则运动上叠加了极小的统计平均偏移。

大截面上各点的流动可能不同，正如河道中央和岸边水速不同。为了描述局部电荷流，引入电流密度
\begin{equation}
  \boxed{\vect j=nq\vect v_d}.
\end{equation}
多种载流子同时存在时
\begin{equation}
  \vect j=\sum_s n_sq_s\vect v_{d,s}.
\end{equation}
穿过任意有向曲面的总电流是电流密度的通量：
\begin{equation}
  \boxed{I=\int_S\vect j\cdot\dd\vect S}.
\end{equation}

\currentDensityDiagram

局部电荷守恒写成连续性方程
\begin{equation}
  \boxed{\frac{\partial\rho_q}{\partial t}+\nabla\cdot\vect j=0}.
\end{equation}
稳恒电流中 $\partial\rho_q/\partial t=0$，所以 $\nabla\cdot\vect j=0$。同一根没有分支的导线中，各截面
电流相同；若流入与流出长期不等，连接处会持续积累电荷，系统就不是稳恒状态。

\subsection*{二端元件与伏安特性}

电路由元件按一定方式连接而成。研究一个元件的外部功能时，可以暂时不追踪内部结构，而只观察端口变量。
最简单的二端元件有两个接线端。先任意规定电压参考正端，并规定参考电流从正端流入；实际方向若与规定相反，
相应物理量取负值。端口极性的选择只是记号，不要与电池自身的正负极混淆。

稳定工作的二端元件在给定端电压 $U$ 时若有唯一电流 $I$，关系 $I=I(U)$ 称为伏安关系，其图像称伏安
特性曲线。不同元件可以有完全不同的曲线，电路设计者只要知道端口特性和连接方式，就能分析整个电路；
器件制造者则还需要研究内部材料与结构。

\ivCharacteristicDiagram

用单值伏安特性描述元件还隐含两个条件。第一，元件附近不能有不可忽略的时变磁通，否则沿不同路径定义的
电压会因电磁感应而不同。第二，元件内部不能持续积累净电荷，否则流入和流出的电流不同，单个端口电流不足
以描述状态。对于具有记忆效应、储能效应或明显暂态的元件，还需要增加内部状态变量，不能只靠静态 $I-U$
曲线概括。

采用被动符号约定时，元件吸收的瞬时功率为
\begin{equation}
  \boxed{P=UI}.
\end{equation}
$P>0$ 表示元件吸收能量，$P<0$ 表示向外部供能。这个符号规则对电阻、电源及一般二端元件都适用。

\subsection*{Ohm 定律：宏观形式与局部形式}

Ohm 定律是导体伏安关系的唯象规律。在温度、材料状态和几何尺寸保持不变时，若电流与电压成正比，定义
比例系数 $R$ 为电阻：
\begin{equation}
  \boxed{U=IR},\qquad I=\frac UR.
\end{equation}
Ohm 定律并不直接说明材料为什么导电，但它给出了工程上极其有效的外部规律，也为建立微观模型提供了实验
约束。只有伏安曲线是过原点直线的元件，才能用与工作点无关的常量 $R$ 描述。

\ohmHistoryContextDiagram

均匀圆柱形导体长度为 $l$、截面积为 $A$，在材料和温度给定时
\begin{equation}
  \boxed{R=\rho\frac lA},\qquad \sigma=\frac1\rho,
\end{equation}
其中 $\rho$ 是电阻率，$\sigma$ 是电导率。长度越大，载流子通过材料时经历的散射过程越多，电阻越大；
截面积越大，可以并行传输电荷的通道越多，电阻越小。

\uniformResistorDiagram

对均匀圆柱导体有 $E=U/l$、$j=I/A$。把几何关系代入 $U=IR$，得到 Ohm 定律的局部形式
\begin{equation}
  \boxed{\vect j=\sigma\vect E}.
\end{equation}
局部形式不显含导体的长度和截面积，更直接体现材料性质。电阻消耗的功率为
\begin{equation}
  \boxed{P=UI=I^2R=\frac{U^2}{R}},
  \qquad
  \boxed{p_V=\vect j\cdot\vect E=\sigma E^2}.
\end{equation}

\subsection*{导电的简化微观模型}

最简单的经典模型把载流子在材料中的阻碍作用类比为与漂移速度成正比的黏滞阻力
\begin{equation}
  \vect f_v=-\alpha\vect v_d.
\end{equation}
稳态时电场力与阻力平衡：$q\vect E-\alpha\vect v_d=0$。因此
\begin{equation}
  \vect j=nq\vect v_d=\frac{nq^2}{\alpha}\vect E,
  \qquad
  \sigma=\frac{nq^2}{\alpha}.
\end{equation}
这个模型极其粗糙，却说明只要平均阻力与速度成正比，就能得到线性 Ohm 定律。

只写 $\vect j=\sigma\vect E$ 会给人一种电场改变后电流瞬间改变的印象。载流子有质量，平均速度的建立必须
经历加速过程。用弛豫时间 $\tau$ 代替阻力系数，Drude 方程写成
\begin{equation}
  m\frac{\dd\vect v_d}{\dd t}
  =q\vect E-\frac m\tau\vect v_d.
\end{equation}
恒定电场在 $t=0$ 突然接通，若初始漂移速度为零，则
\begin{equation}
  \vect v_d(t)=\frac{q\tau}{m}\vect E\left(1-\e^{-t/\tau}\right),
\end{equation}
稳态电导率为
\begin{equation}
  \boxed{\sigma=\frac{nq^2\tau}{m}}.
\end{equation}
电流经过若干个 $\tau$ 才接近 Ohm 定律给出的稳态值；撤去电场后，漂移速度按 $\e^{-t/\tau}$ 衰减，
定向动量在碰撞中传给晶格并转化为热。

\subsection*{电子的无规则运动、漂移与扩散}

金属电子的真实运动并不是沿电场平滑加速再被连续阻力减速。电子以很高速度运动，并不断受到晶格、杂质和
缺陷的散射；外电场只使速度分布产生极小的不对称。用平均自由时间 $\tau$ 描述两次有效散射之间的时间，
电子在这段时间中获得平均漂移速度，碰撞后定向动量被重新随机化。

\drudeDiagram

若平均自由程为 $\ell$、无规则运动典型速率为 $v_t$，则 $\tau\sim\ell/v_t$。在弱场下，电场在一次自由
飞行中造成的速度增量远小于 $v_t$，平均自由程和碰撞率几乎不受外场影响，$\tau$ 可近似视为常量，于是
$j\propto E$。这正是 Ohm 定律成立的重要条件。

载流子除受电场发生漂移外，还会因浓度不均匀从高浓度区域向低浓度区域扩散。若粒子数通量满足
\begin{equation}
  \vect\Gamma=-D\nabla n,
\end{equation}
相应扩散电流密度为 $q\vect\Gamma$。金属稳恒导电中漂移项通常最受关注，而半导体、电解质、气体放电和
结区问题中，漂移与扩散往往必须同时保留。

经典模型还能定性解释电阻的温度依赖。金属升温时晶格振动增强，散射更频繁，$\tau$ 变小；金属载流子浓度
在普通温区变化不大，因此 $\sigma=ne^2\tau/m$ 减小，电阻率增大。半导体升温时虽然散射也增强，但载流子
浓度可能急剧增加，这一效应常占主导，使电阻随温度升高反而下降。两种材料的温度行为不同，根源不在公式
形式，而在 $n(T)$ 与 $\tau(T)$ 的竞争。

以银为例，取 $\sigma\approx6.3\times10^7\,\si{S.m^{-1}}$ 和前述电子数密度，经典估计给出
\begin{equation}
  \tau=\frac{m\sigma}{ne^2}\sim3.8\times10^{-14}\,\si{s}.
\end{equation}
若电子典型速率约为 $10^6\,\si{m.s^{-1}}$，平均自由程
\begin{equation}
  \ell\sim v_t\tau\sim3.8\times10^{-8}\,\si{m},
\end{equation}
比原子间距大数百倍。这个结果说明“电子每遇到一个原子就碰撞一次”的直观图像不正确，也提示必须用波动性、
能带与散射理论建立更精确的金属导电模型。

\subsection*{Ohm 定律的适用范围与气体放电}

得到 $\vect j=\sigma\vect E$ 时使用了若干假设：一次自由飞行中由外场造成的速度增量远小于无规则速度；
碰撞时间不随场强显著改变；散射足够频繁；载流子浓度不随外场改变；材料温度和结构保持稳定。任何条件被
破坏，伏安关系都可能偏离直线。

稀薄气体中的强场导电是典型反例。电子在两次碰撞之间能从电场获得较大能量，碰撞可使分子激发或电离，
载流子浓度随电流改变；放电还可能形成暗放电、辉光放电和弧光放电等不同状态。某些区间甚至出现
$\dd U/\dd I<0$ 的负微分电阻，完全不能用一个常量 $R$ 概括。

\gasDischargeContextDiagram

气体放电曲线的横轴跨越许多数量级。暗放电区电流很小；进入辉光放电后，载流子倍增和空间电荷改变场分布，
端电压可能随电流增加而下降；弧光放电区电流很大而维持电压较低。曲线展示的是材料状态随工作点改变的结果，
不是简单的“电阻数值变了”就能完整解释。

\subsection*{简单电阻网络}

串联电阻中通过各电阻的电流相同，总电压为各电阻电压之和：
\begin{equation}
  U=\sum_iU_i=I\sum_iR_i,
  \qquad
  \boxed{R_{\rm series}=\sum_iR_i}.
\end{equation}
并联电阻两端电压相同，总电流为各支路电流之和：
\begin{equation}
  I=\sum_i\frac U{R_i},
  \qquad
  \boxed{\frac1{R_{\parallel}}=\sum_i\frac1{R_i}}.
\end{equation}

\resistorConnectionDiagram

能逐级化成串、并联的电路称为简单网络。很多实际电路含有桥式连接、多个电源或非线性元件，无法只靠串并联
等效化简，这时必须使用一般电路网络定律。

\subsection*{电路网络术语}

电路由元件及连接线组成。两个或更多元件的连接点称为节点，相邻节点之间的一段通路称为支路；从某个节点
出发，沿若干支路行进并最终回到出发节点，构成一个回路。只有两条支路相连的节点通常不会产生新的电流分配
条件；三条或更多支路相连的非平凡节点才需要独立列 KCL。随着网络规模增大，可画出的闭合回路很多，但它们
并不全部独立。

\circuitTerminologyDiagram

\subsection*{Kirchhoff 电流定律与电压定律}

Kirchhoff 电流定律（KCL）是电荷守恒在集总电路中的表现：任一节点流入电流的代数和为零，等价地，流入
节点的电流总和等于流出总和：
\begin{equation}
  \boxed{\sum_k I_k=0}.
\end{equation}
列式前可以任意假设各支路电流方向。求得某电流为负，只说明真实方向与假设相反，并不表示方程错误。含有
$N$ 个节点的连通网络只有 $N-1$ 个独立 KCL 方程，因为把全部节点方程相加会得到恒等式。

Kirchhoff 电压定律（KVL）指出，在没有不可忽略时变磁通穿过回路的集总电路中，沿任一闭合回路各支路电压
的代数和为零：
\begin{equation}
  \boxed{\sum_k U_k=0}.
\end{equation}
选定绕行方向后，从元件负端走向正端记为电势升高，反向记为电势降低；只要全程符号规则一致即可。KVL 可
看成静电势单值性或稳恒电路能量守恒的表达。如果回路穿过时变磁通，应使用 Faraday 定律
$\sum U_k=-\dd\Phi_B/\dd t$，不能继续把右侧写成零。

\kirchhoffDiagram

回路方程也有线性相关性。选取一组能覆盖网络中全部独立闭合路径的基本回路即可；多列出的回路通常是已有
方程的线性组合。KCL、KVL 再配合每个元件的伏安关系，原则上能够确定全部支路电流和元件电压。

\subsection*{节点电压法}

直接把每条支路电流和电压都设为未知量会产生很大的方程组。节点电压法先任选一个节点为参考地，令其电势
为零，再把其余节点电势作为未知量。支路电压自动等于两端节点电势之差，KVL 因而自动满足；对每个未知电势
节点列 KCL，再用元件伏安关系把支路电流写成节点电势，即可得到较小的方程组。

\nodeMethodDiagram

例如图中参考节点为零电势，理想电源令节点 $a$ 的电势为 $\mathcal E_0$，只有节点 $b$ 的电势 $\varphi$
未知。对 $b$ 列 KCL 时，每条电阻支路电流均写成
$(\varphi-\varphi_{\rm neighbor})/R$，所得只是关于 $\varphi$ 的一次方程。解出节点电势后，再逐支路求
电压和电流。经验上把连接较多支路或电压源的节点选为参考地，往往能减少未知量和代数步骤。

\subsection*{电流图像中容易混淆的三个时间尺度}

导体与绝缘体的区别在于，导体内部存在能够在外场作用下产生宏观迁移的载流子。金属、电解质、半导体、
电离气体和超导体虽然都能导电，内部机制却截然不同。金属中有近似离域的电子；电解质中正负离子向相反
方向漂移却可能产生同向的传统电流；半导体的载流子浓度会受温度、掺杂和光照显著影响；气体通常要先被
电离才有足够载流子；超导态更不能用普通散射模型解释。把这些材料都写成一个电阻符号，是电路层次的抽象，
并不意味着它们在微观层次遵循同一种过程。

经典电子模型的局限需要在开始时说明。金属电子不能用经典能量均分定理确定速度，它们服从量子统计；周期
晶格中的运动也必须用能带描述。这里使用漂移速度、平均自由时间和有效碰撞，是为了提取与宏观电流直接相关
的统计量。模型只要能够解释 $j\propto E$、Joule 热、暂态响应和温度趋势，就已经完成本讲的任务；它不能
替代近代固体物理理论。

银导线中电子的漂移速度只有 $10^{-4}\,\si{m.s^{-1}}$ 的数量级，不能据此认为合上开关后，要等电子从电源
缓慢漂移到灯泡，灯才会亮。这里存在三个完全不同的时间尺度。第一是电磁场和表面电荷重新分布的传播时间，
信号沿电路以接近电磁波在周围介质中的传播速度建立；第二是电子平均漂移速度发生改变所需的弛豫时间，典型
金属中可短到 $10^{-14}\,\si{s}$；第三才是某个电子在很长导线中产生可观净位移所需的漂移时间。电路快速
响应由前两个尺度决定，不是靠同一批电子从电源跑到负载。

电流定义中的“通过截面”也必须带方向。若正电荷沿截面法向通过，记为正电流；负电荷反向通过同样产生正
电流。对于电子导电，电子速度与传统电流方向相反，但 $q<0$，所以 $\vect j=nq\vect v_d$ 仍自动给出正确
方向。多载流子系统中不能简单比较粒子速度，还要把每类载流子的电荷符号与浓度一起计入。

稳恒电流并不要求每个电子做匀速直线运动，也不要求任一点没有电荷。它只要求宏观电流分布不随时间改变，
局部电荷密度也不再持续积累。导线表面可以存在极小而稳定的电荷梯度，正是这些表面电荷建立了导线内部所需
的电场。若在某节点流入电流大于流出电流，节点电荷会增加并改变周围电势，直到新的电场把各支路电流调整到
满足 KCL 的状态。

\subsection*{从端口观点理解元件}

把元件画成方框并不是说内部结构不重要，而是根据研究目的选择描述层次。分析收音机电路时，可能只需要晶体
管的端口特性；设计晶体管时，则必须研究载流子的注入、复合和能带。端口模型的价值在于模块化：只要每个元件
的外部关系和连接拓扑已知，就可以不重复求解内部所有自由度。

规定端口电压、电流方向后，伏安图必须明确坐标轴。若横轴为 $U$、纵轴为 $I$，过原点直线的斜率是
$I/U=1/R$；若交换坐标轴，斜率才是 $R$。同一工作点处还可以区分静态电阻 $U/I$ 与微分电阻
$\dd U/\dd I$。非线性元件的这两个量通常不同，负微分电阻表示局部斜率为负，并不等于元件在该点的
$U/I$ 必然为负。

理想电阻的伏安曲线经过原点，正、负电压下具有反向对称性。实际金属电阻通电后会发热，温度改变又使电阻率
改变，若测量过程不能保持温度恒定，伏安曲线也会弯曲。因此，Ohm 定律中的“外部状态不变”不是可有可无的
附加语，而是把 $R$ 当作常量的必要条件。

元件功率的正负由参考方向共同决定。采用被动符号约定时，电流从电压正端流入，$P=UI>0$ 表示电能进入元件；
电源正常供能时，实际电流常从其正极流出，相对于被动约定 $I<0$，因此 $P<0$。只要参考方向定义清楚，
不需要为电阻和电源分别发明两套功率公式。

静态伏安曲线还要求测量过程足够慢，使元件内部状态已经稳定。电容器在直流稳态下电流为零，却在充电时有
电流；电感在直流稳态下近似导线，却在电流变化时产生感应电压。若只给出某一瞬间的 $U,I$，不能判断这些
储能元件的后续行为。一般电路理论之所以要把电阻、电容、电感和受控源分开，就是因为它们需要不同的状态
关系。当前章节先集中研究不储存可恢复电磁能的电阻型元件。

\subsection*{从宏观电阻到微观电导率}

$R=\rho l/A$ 把材料性质与几何性质分开。将同一种均匀电阻丝长度加倍，相当于把两个相同电阻串联，电阻
加倍；将两根相同电阻丝并排连接，相当于截面积加倍，电阻减半。这一几何推理与连续介质计算完全一致。
电阻率 $\rho$ 才是给定温度和材料状态下更接近材料本身的物理量。

局部 Ohm 定律可以从圆柱导体反推。均匀导体内 $E=U/l$、$j=I/A$，代入
$U=I\rho l/A$，消去 $l,A$ 得 $E=\rho j$。各向同性材料中 $\vect j$ 与 $\vect E$ 同向，写成
$\vect j=\sigma\vect E$。若材料各向异性，电导率应当是张量，电流密度与电场不一定平行；若材料不均匀，
$\sigma$ 还会随位置变化。

体积功率密度 $p_V=\vect j\cdot\vect E$ 可由单个电荷的功率相加得到。单位体积内第 $s$ 类载流子的受力
功率为 $n_sq_s\vect v_{d,s}\cdot\vect E$，对全部载流子求和正是 $\vect j\cdot\vect E$。在线性被动介质中
$p_V=\sigma E^2\geq0$，这些能量最终通过散射传给晶格，表现为 Joule 热。

\subsection*{两种经典导电模型的联系}

黏滞阻力模型写作 $m\dot{\vect v}_d=q\vect E-\alpha\vect v_d$。与 Drude 形式比较可知
$\alpha=m/\tau$，两种写法只是参数选择不同。恒定电场下方程的完整解为
\begin{equation}
 \vect v_d(t)=\vect v_{\infty}+
 \left[\vect v_d(0)-\vect v_{\infty}\right]\e^{-t/\tau},
 \qquad
 \vect v_{\infty}=\frac{q\tau}{m}\vect E.
\end{equation}
它同时包含电流建立和撤场后的衰减。若外场变化周期远大于 $\tau$，电子几乎总能跟上瞬时稳态，静态 Ohm
定律是好近似；若频率高到与 $1/\tau$ 同量级，就必须保留惯性项，电导率变成与频率有关的复数量。

用平均自由飞行图像估计漂移速度时，电子在一次碰撞后若完全失去定向动量，随后在电场中从零加速。一次自由
飞行中的平均定向速度严格含有约 $1/2$ 的因子；但碰撞时间有统计分布，简单数量级估计常把这个系数吸收到
$\tau$ 的定义中。竞赛推导应先说明所用模型，再保持同一套定义，避免一处用平均时间、一处用固定碰撞时间
而造成因子不一致。

经典模型把碰撞画成电子与固定原子小球相撞，这只是示意。周期晶格本身并不会像随机障碍物那样阻止理想 Bloch
电子运动，真正产生电阻的是晶格振动、杂质、缺陷和边界等破坏平移对称性的散射机制。银的平均自由程远大于
原子间距，正是朴素“小球连续碰撞”图像失效的证据之一。

\subsection*{温度、场强和载流子浓度}

电导率 $\sigma=nq^2\tau/m$ 中的每个量都可能依赖外部条件。普通金属的导电电子浓度主要由价电子数和原子
密度决定，在不发生相变的温区变化较小；升温主要增强声子散射，使 $\tau$ 缩短，电阻率上升。常温附近常用
近似
\begin{equation}
 R(T)=R_0\left[1+\alpha_R(T-T_0)\right]
\end{equation}
描述，但温区很宽时必须使用非线性关系。

本征半导体中的载流子需要由热激发跨越能隙产生，$n(T)$ 对温度极为敏感。升温虽然也会缩短散射时间，却可能
使载流子数量增加得更快，因而总电导率上升。掺杂半导体还会经历杂质电离、饱和和本征激发等不同温区，不能
用一句“半导体电阻随温度下降”覆盖全部情况。

强电场会使电子在一次自由飞行中获得不可忽略的能量，改变碰撞截面、迁移率甚至载流子浓度。气体中电子可
通过碰撞电离产生更多电子和离子，形成雪崩；半导体中也可能出现速度饱和、碰撞电离和击穿。此时 $n$ 与
$\tau$ 都随场改变，线性 Ohm 定律自然失效。

气体放电曲线上不同分支对应不同放电机制和空间结构。同一电压可能对应不止一个稳定或暂稳电流，元件状态还
取决于过去的放电历史。分析这类元件时，不能仅用一个静态电阻值替代完整伏安曲线，更不能把曲线某段的斜率
直接当成始终有效的电阻。

Ohm 定律是唯象定律而不是自然界无条件成立的基本定律。它的地位类似小形变下的 Hooke 定律：在适用范围内
非常准确，超出范围后必须改用更一般的材料关系。工程师频繁使用 Ohm 定律，恰恰是因为许多金属元件在正常
温度与电流范围内足够接近线性；这并不能推出所有导电现象都应当线性。判断题目能否使用 $U=IR$ 时，应先
检查材料、温度、场强和工作状态是否允许把 $R$ 当作常量。

\subsection*{网络方程为什么会有冗余}

对含 $N$ 个节点的连通网络逐节点列 KCL，会得到 $N$ 个方程；但每一支路电流在两个端点方程中分别以相反
符号出现，把全部节点方程相加后左边恒为零。因此独立 KCL 只有 $N-1$ 个。少掉的一个方程并不是信息丢失，
而是整个网络总电荷守恒已经包含在其余节点方程中。

KVL 的冗余可用回路叠加理解。如果回路 $L_3$ 的有向边恰好等于回路 $L_1$ 与 $L_2$ 的有向边之和，那么
$L_3$ 的电压方程就是前两个方程相加。对于有 $B$ 条支路、$N$ 个节点的连通网络，独立回路数为
$B-N+1$。选择一棵包含全部节点而没有回路的生成树，每增加一条树外支路就产生一个独立基本回路。

KCL 给出 $N-1$ 个独立约束，KVL 给出 $B-N+1$ 个独立约束，两者合计 $B$ 个，再配合 $B$ 条支路各自的
伏安关系，恰好构成求解全部支路电流和电压的完备方程组。这一计数解释了为什么既不能少列必要方程，也没有
必要把图中所有能看到的回路都列一遍。

KCL 与 KVL 的适用也包含集总近似。电路尺寸远小于信号变化对应的电磁波波长时，可以把导线近似为瞬时等势
连接，把电场和磁场效应集中到元件参数中。若电路很大或频率很高，电压、电流会沿传输线传播，同一根导线的
不同位置不再处于同一瞬间状态，简单节点和支路模型需要升级为分布参数模型。竞赛中的直流网络通常默认满足
集总条件，但理解这一前提能防止把 KVL 错用于穿过显著时变磁通的回路。

\subsection*{节点电压法的列式模板}

节点法处理纯电阻支路时，可固定参考节点电势为零。若未知节点 $i$ 与相邻节点 $j$ 通过电阻 $R_{ij}$ 相连，
从 $i$ 流向 $j$ 的电流写成
\begin{equation}
 I_{ij}=\frac{\varphi_i-\varphi_j}{R_{ij}}.
\end{equation}
对节点 $i$ 列 KCL：
\begin{equation}
 \sum_j\frac{\varphi_i-\varphi_j}{R_{ij}}=I_i^{\rm inj},
\end{equation}
其中右侧表示外部电流源注入该节点的净电流。所有未知节点都得到一个线性方程，系数矩阵的对角元是与该节点
相连的电导之和，非对角元是相邻节点间电导的负值。

若理想电压源连接未知节点与参考地，该节点电势直接由电源确定；若理想电压源跨接两个未知节点，可把两节点
组成超节点，对超节点整体列 KCL，再增加电压约束。节点法并没有废除 KVL，而是用节点电势差自动保证所有
不含时变磁通的回路满足 KVL。


\subsection*{重点知识}

\begin{enumerate}[label=\textbf{33.\arabic*},leftmargin=4.2em]
  \item 电流是电荷通过截面的速率，$I=\dd Q/\dd t$；电流密度满足
  $I=\int_S\vect j\cdot\dd\vect S$。
  \item 二端元件的伏安特性描述端电压与端电流的关系；采用被动符号约定时功率 $P=UI$。
  \item Ohm 元件满足 $U=IR$，均匀导体 $R=\rho l/A$，局部形式为 $\vect j=\sigma\vect E$。
  \item Drude 模型给出 $\sigma=nq^2\tau/m$，同时说明电流建立存在弛豫时间，Ohm 定律有适用条件。
  \item KCL 来自电荷守恒，KVL 来自电势单值性与能量守恒；节点电压法可系统求解复杂网络。
\end{enumerate}
